% ============================================================
% Section 6.1 & 6.4: New Discussion sections for WISE
% Quality Assurance for Web Systems and Deployment Implications
% ============================================================

\subsection{Quality Assurance for Trustworthy Web Systems}

Web information systems serving critical applications---medical diagnosis support,
legal research, financial advisory, scientific literature review---require not
only accurate answers but \emph{trustworthy} answers: outputs whose quality can
be assessed, whose reasoning can be traced, and whose failures can be detected
and explained. The keynote theme for WISE 2026, ``Ensuring Quality in LLM-Driven
Agents''~\cite{benatallah2026keynote}, emphasizes this shift from pure accuracy
metrics to comprehensive quality assurance as the central challenge for deploying
LLM-based Web systems in production.

\sys demonstrates three architectural principles for trustworthy GraphRAG that
directly address this challenge:

\smallskip
\noindent\textbf{1. Verification as first-class component, not afterthought.}
The VA is not a post-processing check or an optional evaluation step---it is a
core component of the inference pipeline, executed on every question, with its
output directly steering the retrieval-repair loop. This design ensures that
\emph{every} answer comes with a quality signal ($s \in [0,1]$) and a diagnostic
record, making quality assessment intrinsic rather than external. The 24.2\%
convergence-correctness gap (Section~\ref{sec:quality_analysis}) validates that
this intrinsic verification provides actionable quality information.

\smallskip
\noindent\textbf{2. Traceability by design, not logging.}
The VA's diagnostic output (Eq.~\ref{eq:feedback}) is not an execution log or
attention heatmap but a \emph{structured natural-language record} identifying
which sub-query failed, why, and what specific information is missing. This
record serves dual purposes: it enables surgical repair (Figure~\ref{fig:case_study_flow})
and it provides auditability for regulated domains. A medical QA system can
present the VA's diagnostic to a clinician explaining why an answer is uncertain;
a legal research tool can log the diagnostic for compliance auditing; a
conversational agent can surface it as a user-facing explanation of retrieval
limitations.

\smallskip
\noindent\textbf{3. Bounded quality guarantees for production deployment.}
Proposition~1 bounds inference cost at $4N$ LLM calls; Proposition~2 guarantees
monotone best-answer tracking. These formal properties ensure that \sys's quality
assurance mechanism comes with \emph{predictable overhead} and a \emph{reliable
quality floor}---essential for Web systems operating under latency constraints
or serving queries at scale. The average of 1.62 iterations (Table~\ref{tab:iteration_breakdown})
confirms that the bound is tight in practice.

The quality analysis in Section~\ref{sec:quality_analysis} validates these
principles empirically: VA convergence predicts correctness (24.2\% gap), VA
scores discriminate answer quality (0.126 separation), and surgical repair
recovers from retrieval failures (61.8\% questions benefit from refinement).
Together, these findings establish that sub-query-level quality verification
is not merely a performance optimization but a fundamental architectural shift
toward trustworthy GraphRAG for Web systems.

The broader implication for the WISE community: as GraphRAG becomes the dominant
paradigm for knowledge-intensive Web applications, the research challenge shifts
from \emph{retrieval accuracy} (can we retrieve the right triples?) to
\emph{trustworthy retrieval} (can we verify that retrieval succeeded, explain
when it failed, and provide quality guarantees for production deployment?). \sys
demonstrates that this shift is architecturally feasible through verification-driven
design, and the empirical gains validate that trustworthiness and accuracy are
not in conflict---they reinforce each other.

\subsection{Implications for Web System Deployment}

Deploying GraphRAG systems in production Web environments requires careful
consideration of the accuracy--cost--trustworthiness tradeoff. We discuss \sys's
deployment characteristics and compare them to the strongest baselines from the
1000-question HotpotQA evaluation (Table~\ref{tab:hotpot1000}).

\smallskip
\noindent\textbf{Inference cost.}
\sys averages 1.62 iterations $\times$ 4 agents = 6.5 LLM calls per question,
plus $\sim$5 calls for task-specific KG construction, totaling $\sim$11.5 calls
per question. This is $\sim$11.5$\times$ Naive RAG (1 call) but comparable to
other multi-step systems: A-RAG averages 5.0 calls (4.03 retrieval loops + 1
final answer extraction), IRCoT averages 2.75 calls. HopRAG does not report
per-question call counts but performs retrieve-reason-prune over a passage graph,
implying multiple LLM invocations per question.

The cost premium is justified by the dual benefits: \sys delivers both
\emph{competitive accuracy} (55.3\% EM pure selector, 58.4\% EM with consensus
reranking; Table~\ref{tab:hotpot1000}) and \emph{trustworthy outputs} (quality
scores, convergence signals, traceable diagnostics). For Web applications where
answer trustworthiness is critical---medical QA, legal research, financial
advisory---this tradeoff is favorable.

\smallskip
\noindent\textbf{When to use \sys vs. alternatives.}
\begin{itemize}
  \item \textbf{Use \sys when}: (1)~the application domain requires auditable
        reasoning (healthcare, finance, legal), (2)~silent retrieval failures
        are unacceptable (safety-critical systems), (3)~downstream systems need
        per-answer confidence signals for decision-making, or (4)~the system
        must explain \emph{why} it cannot answer rather than returning a
        likely-wrong answer.

  \item \textbf{Use A-RAG or HopRAG when}: Raw accuracy is the primary objective,
        latency constraints are tight, and quality assurance is handled externally
        (e.g., human review of all outputs). A-RAG reaches 51.7\% EM with adaptive
        multi-step retrieval; HopRAG reaches 54.7\% EM with passage-graph pruning.
        Both are faster than \sys but provide no intrinsic quality verification.

  \item \textbf{Use Naive RAG or LightRAG when}: The task does not require
        multi-hop reasoning, the KG construction overhead is prohibitive, or the
        system operates in a high-throughput, low-latency regime where single-pass
        retrieval is mandatory.
\end{itemize}

\smallskip
\noindent\textbf{Latency and throughput considerations.}
At $\sim$11.5 LLM calls per question, \sys's end-to-end latency depends on the
LLM backend. With a frontier model (e.g., GPT-4, Claude) at $\sim$2--5s per call,
total latency is $\sim$23--58s per question assuming sequential execution. This
is acceptable for non-interactive Web applications (research tools, report
generation, batch QA) but too slow for conversational agents requiring
sub-second response times. Two mitigation strategies:
\begin{enumerate}
  \item \textbf{Parallel agent execution.} QDA, GRA, AGA calls within a single
        iteration can be pipelined (GRA waits for QDA; AGA waits for GRA), but
        the four agents of iteration~$t$ need not block iteration~$t{+}1$ if
        speculative execution is permitted. A speculative second iteration can
        begin before the VA completes, trading compute for latency.

  \item \textbf{Hybrid tiering.} Route simple questions (single-hop, high
        initial BM25 confidence) to Naive RAG; route complex questions requiring
        multi-hop reasoning or low-confidence initial retrieval to \sys. A
        lightweight classifier (e.g., question type, BM25 top-score gap) can
        perform this routing with negligible overhead.
\end{enumerate}

\smallskip
\noindent\textbf{Persistent vs. task-specific KGs.}
\sys builds a task-specific KG per question from supporting passages
(Definition~4), averaging 148 nodes and 213 edges per HotpotQA question. This
trades KG construction cost ($\sim$5 LLM calls) for compactness (no irrelevant
triples) and eliminates large-scale entity disambiguation. For enterprise
deployment, a \emph{persistent domain KG} built offline from a curated corpus
eliminates per-question KG construction but requires (1)~a disambiguation layer
for entity linking and (2)~KG maintenance as the knowledge source evolves. The
modular design (Design Principle~P1, Section~\ref{sec:principles}) ensures that
swapping task-specific KG construction for persistent KG lookup requires only
GRA prompt adjustments with no change to the VA, QDA, or AGA.

\smallskip
\noindent\textbf{Regulatory and compliance contexts.}
For Web systems subject to regulatory oversight (e.g., FDA-regulated clinical
decision support, financial advisory under fiduciary duty, legal research tools
in discovery), \sys's traceable diagnostics provide a compliance advantage over
black-box GraphRAG systems. The VA's diagnostic record can be logged for audit
trails, presented to domain experts for validation, or surfaced to end-users as
explanations. The 24.2\% convergence-correctness gap (Table~\ref{tab:convergence_quality})
provides a quantitative basis for setting quality thresholds: a system requiring
95\% precision might refuse to answer when $s < 0.7$, trading recall for
trustworthiness.

\smallskip
\noindent\textbf{Generalization beyond multi-hop QA.}
The three design principles (Section~\ref{sec:principles}) apply to any Web
knowledge retrieval task where intermediate steps can fail silently:
entity-centric fact verification (verify claims against a KG, detect unsupported
statements), structured Web search (retrieve entity attributes, verify
completeness), open-domain QA with KG augmentation (retrieve from both text and
KG, verify consistency). The verification-driven architecture is task-agnostic;
only the QDA and GRA prompts need task-specific adaptation.

\smallskip
\noindent\textbf{Summary.}
\sys is production-ready for Web applications prioritizing trustworthiness over
raw throughput. Its bounded cost ($\sim$11.5 calls/question), monotone quality
guarantees (Proposition~2), and training-free deployment (no fine-tuning, no
proprietary data) make it deployable via API-accessible LLMs. The quality
assurance mechanisms---convergence signals, diagnostic records, sub-query-level
verification---address the core requirements for regulated domains and
safety-critical Web systems. For applications where silent retrieval failure is
unacceptable, \sys provides an architecturally sound path to trustworthy GraphRAG.

